\section{实验设置与评价协议}
\subsection{受试者分组与模型选择}
仅在官方训练集内使用五折 GroupKFold，将同一受试者的全部窗口放入同一折。五个外层验证组依次包含受试者 $\{14,15,19,25\}$、$\{6,16,21,22\}$、$\{3,11,17,26\}$、$\{1,7,23,30\}$ 和 $\{5,8,27,28,29\}$，对应验证窗口数为 1420、1420、1417、1410、1685。六类模型使用相同外层划分，避免重叠窗口跨训练与验证集合共享同一人的信息。

传统模型按表\ref{tab:grid}的网格搜索，以五折宏平均 F1 选择参数，逻辑回归最大迭代次数为 3000，随机森林随机种子为 42。标准化器放在分类流水线内，在每折训练部分拟合。这里的五折分数同时用于选择传统模型超参数与报告开发结果，存在选择乐观性，并非完全嵌套交叉验证的无偏泛化估计。
\begin{table}[tbp]\centering\small
\caption{传统分类器搜索空间与选定参数}\label{tab:grid}
\begin{tabularx}{\linewidth}{lYY}\toprule
模型 & 搜索空间 & 选定参数\\\midrule
逻辑回归 & $C\in\{0.001,0.01,0.1,1,10\}$ & $C=0.1$\\
RBF SVM & $C\in\{1,10,100\}$；$\gamma\in\{\texttt{scale},0.001,0.01\}$ & $C=10,\ \gamma=0.001$\\
随机森林 & 深度 $\in\{\text{不限},20\}$；叶节点最小样本数 $\in\{1,3\}$ & 深度 20；最小样本数 3\\\bottomrule
\end{tabularx}
\end{table}

神经网络结构保持固定，每个外层训练集再用 GroupShuffleSplit 划出约 20\% 的受试者作为内层验证组，划分种子固定为 42。此比例指受试者比例，不要求窗口数恰为 20\%。首先仅用内层训练部分拟合预处理参数并训练，按内层宏平均 F1 选择轮数。上限为 100 轮，连续 10 轮无严格提升即停止。随后重新初始化网络，在全部外层训练数据上重新拟合标准化参数，并按选定轮数重训，最后仅评价一次外层验证集。外层验证分数不参与该折的早停。

三类神经网络均使用批量大小 128、学习率 $10^{-3}$、权重衰减 $10^{-4}$ 和交叉熵损失。基础种子为 42 时，第 $k$ 折训练种子为 $42+k-1$。融合模型另外使用基础种子 2024、2025 复验，保持数据划分、结构和优化设置一致，仅改变初始化及训练随机性。推断时使用评价模式，关闭 Dropout，并使用 BN 的运行统计量。

\subsection{最终训练与独立测试}
根据训练集开发结果，在查看官方测试标签对应的预测成绩前，将逻辑回归 $C=0.1$、不使用 PCA 固定为主基准，融合模型固定为比较模型。融合最终轮数取三个基础种子共 15 个内层所选轮数的中位数 9，最终训练种子固定为 42。两模型用完整的 7352 个训练窗口重新训练并保存模型、预处理参数及配置，随后加载这些文件对 2947 个官方测试窗口进行预测。测试时不重新拟合标准化器或更新模型权重。本次仅对这两个冻结模型实施独立测试。

\subsection{指标与实现环境}
以准确率和宏平均 F1 为主要指标，六个类别等权。令 $TP_c,FP_c,FN_c$ 分别表示类别 $c$ 的真阳性、假阳性和假阴性数量，则
\begin{equation}
 \mathrm{Accuracy}=\frac{\sum_i\mathbf1(\widehat y_i=y_i)}{N},\quad
 F_{1,c}=\frac{2TP_c}{2TP_c+FP_c+FN_c},\quad
 F_{1,\mathrm{macro}}=\frac16\sum_{c=1}^{6}F_{1,c}.
\end{equation}
分母为零时按实现约定将相应 F1 记为零。五折分数按折等权求均值，折间标准差使用总体形式（ddof=0），不能视为独立测试指标或置信区间。三次复验则对三个五折均值另算样本标准差（ddof=1），与折间波动分别报告。

实验在 Windows 的 Python 3.14.3 环境执行，NumPy、scikit-learn 与 PyTorch 版本分别为 2.5.4、1.9.1 和 2.14.1+cpu，采用 CPU 计算。神经网络设置四个计算线程，DataLoader 工作进程数为零。传统分类器与深度学习分别基于 scikit-learn\cite{sklearn2011}和 PyTorch\cite{pytorch2019}。各阶段结果以 CSV/JSON、逐样本预测和模型文件留档；绘图与论文表格从这些结果文件生成，不重新训练或预测。文中耗时仅描述本机单次运行，不作为严格控制条件下的计算效率排名。
